\section{地球の層とプレート運動をつなぐ}
\label{theory:interior:overview}
地殻は物質の種類による区分、プレートは長い時間でも比較的硬く一体として動く部分の区分である。まず分類軸を分け、その後に地震波・静力学・熱輸送を同じ地球に適用する。典型的な数値は地域による違いを持つ推定であり、普遍定数として暗記しない。

\subsection{観測できない内部をどのように調べるか}
\label{theory:interior:observations}
内部の多くは直接採取できない。地震波の走時と反射・屈折は弾性と密度、重力は質量分布、地表熱流束は熱輸送を制約する。複数の観測と岩石の高圧実験を合わせて組成・温度を推定する。地震波速度の変化だけで全ての境界を同じ種類の物質境界と断定しない。\cite{usgs_interior_structure}

\subsection{大陸地殻と海洋地殻の典型像}
\label{theory:interior:crust}
海洋地殻は海嶺でのマグマ形成と固化により生まれ、玄武岩と斑れい岩を主体とする。大陸地殻は火成・堆積・変成・変形の長い履歴を持ち、上下にも水平方向にも不均質である。花崗岩層と玄武岩層という古い二層像は平均組成の近似としてのみ用いる。海水を除く海洋地殻が数km、大陸地殻が数十kmという桁の差を先に理解する。\cite{umich_crust_teaching,usgs_lithosphere2002}

\subsection{モホ面は地殻底を表す}
\label{theory:interior:moho}
モホロビチッチ不連続面は地殻とマントルの間に認められる地震学的な境界である。深さは場所により異なり、海面から測るか固体表面から測るかでも数値が変わる。海水厚さを地殻厚さに含めない。地殻底がそのままプレート底であるとはいえない。\cite{usgs_lithosphere2002}

\subsection{リソスフェアとアセノスフェア：時間尺度を含む硬さ}
\label{theory:interior:lithosphere}
リソスフェア（岩石圏）は地殻と最上部マントルを含み、プレートとして比較的まとまって動く。アセノスフェア（岩流圏）はその下で長期的に変形しやすい領域である。「流れる」は液体であることを必ずしも意味せず、主に固体の岩石も長時間には粘性変形する。力学的な境界は採用する観測と時間尺度に関わり、一定深度の一枚の面を全地球に当てはめない。\cite{kagoshima_plate2019}

\begin{center}
\begin{tabular}{lll}
\toprule
区分の軸 & 上側 & 下側 \\
\midrule
組成・地震学的な地殻底 & 地殻 & マントル \\
長期の力学的な変形 & リソスフェア & アセノスフェア \\
\bottomrule
\end{tabular}
\end{center}
上下の行は同じ位置に境界を置いた断面図ではない。これが令和7年度問24の判別に必要な点である。

\subsection{上載荷重から静岩圧へ}
\label{theory:interior:pressure}
水平で静的な上載荷重だけを考え、下向き深さを\(z\)、密度を\(\rho(z)\)、重力加速度を\(g(z)\)とする。薄層の重さを積分すると
\[
 p(z)-p(0)=\int_0^z\rho(s)g(s)\,ds.
\]
密度と重力が一定なら\(p-p(0)=\rho gz\)である。\(\mathrm{kg\,m^{-3}}\cdot\mathrm{m\,s^{-2}}\cdot\mathrm m=\mathrm{Pa}\)であり、地殻密度約\(10^3\)と数十kmの積は\(10^9\)Pa程度となる。平成27年度問32は平均密度・厚さ・重力・間隙流体無視を指定するため、この式を直接用いられる。間隙水圧、有効応力、異方的な応力場を静岩圧の値だけで代替しない。

\subsection{大陸の上下組成と面積の分母}
\label{theory:interior:continental}
大陸上部は平均的に珪長質寄り、下部はより苦鉄質寄りと整理する。岩石名をそのまま均質層の保証とせず、下部の高い密度という傾向と結び付ける。大陸棚の下も大陸地殻であるため、大陸地殻の表面積割合約4割は、海面より上にある陸地の割合約3割と異なる。\cite{umich_crust_teaching}

\subsection{岩石の年代と海洋底の更新}
\label{theory:interior:ages}
大陸の一部には約40億年前の岩石が残る。一方、海洋底は海嶺で生成して移動し、沈み込みにより再循環するため、典型的な現存海洋底は数億年以内である。地球誕生の年齢、鉱物一粒の形成、岩体の固化、変成の時期は別の時計である。大洋底を30億年までとする肢には、この更新の仕組みとの不整合がある。\cite{tokyo_crust_gallery,usgs_plate_boundaries}

\subsection{発散・収束・横ずれ}
\label{theory:interior:boundaries}
発散境界では二つのプレートが離れ、新しい海洋地殻が生じる。収束境界では接近し、沈み込みや衝突が起こる。トランスフォーム境界は相対的な横ずれである。海嶺のずれを結ぶ活動区間と、その外側へ延びる非活動的な断裂帯を分ける。伸張に伴う正断層をトランスフォームの基本運動としない。\cite{usgs_plate_boundaries}

\subsection{海嶺の高さと拡大速度}
\label{theory:interior:ridge}
熱膨張を一次近似すると\(\rho(T)\simeq\rho_0[1-\alpha(T-T_0)]\)である。高温の新しい海洋リソスフェアは周囲の冷えた部分より低密度であり、浮力と海嶺の地形の高まりを結び付けられる。温度・組成・厚さの違いもあるため、この式だけで海嶺の高さを精密に計算したとはしない。
超低速拡大の一部ではマグマ供給が小さく、マントル岩が海底へ露出する。ゆっくり拡大する全区間が火成活動を欠く、又は熱水活動を欠くという意味ではない。\cite{noaa_ultraslow2004}

\subsection{相対速度には参照系と地点が必要である}
\label{theory:interior:velocity}
球面上の剛体プレートの速度は\(\boldsymbol v=\boldsymbol\omega\times\boldsymbol r\)と表せる。二つのプレートの相対角速度は\(\boldsymbol\omega_A-\boldsymbol\omega_B\)であり、同じプレート対でも位置\(\boldsymbol r\)により速度の大きさが変わる。日本付近の代表的な太平洋・フィリピン海プレートの沈み込み速度はそれぞれ年8〜9 cm・4〜5 cmの桁である。令和元年度問30の9 cmと4.5 cmはこの代表比較に対応するが、北米固定という指定と採用モデルの厳密な照合は査読対象である。\cite{kagoshima_plate2019}

\subsection{島弧の位置関係}
\label{theory:interior:arc}
沈み込みの海溝側から、前弧、火山フロント、背弧という位置関係を捉える。海溝から大陸側へ深くなるスラブ内の地震分布が和達・ベニオフ帯である。背弧は収束するプレート境界の上盤側でも伸張して海盆を作る場合がある。縁海という地理用語と背弧海盆という形成過程による用語を同一の定義としない。\cite{usgs_plate_boundaries}

\subsection{温度・地温勾配・熱流量を分ける}
\label{theory:interior:heat}
地表で測る地殻熱流量は単位面積・単位時間当たりの熱輸送であり、\(\mathrm{W\,m^{-2}}\)を持つ。熱伝導では\(\boldsymbol q=-k\nabla T\)である。下向き深さ\(z\)と上向きを正の流束では\(q_{\rm up}=k\,dT/dz\)となる。地下水の移流もある場合、温度勾配だけから総熱輸送を決められない。
東北日本の広域分布は太平洋側で低く、日本海側で高い。令和元年度再試験問33はこの大小を逆転させた肢を含む。1991年の原研究要旨で出題以前の分布を確認しており、現在の一点の測定を過去の全球モデルへ置き換えてはいない。\cite{gsj_heat_flow1991}

\subsection{地殻内地震と活断層の繰返し}
\label{theory:interior:crust_quakes}
地殻内の断層の脆性破壊と、沈み込み境界面のすべりを分ける。震央が沿岸海域にあることだけでは地殻内地震を除外できない。1995年兵庫県南部地震の気象庁M7.3は、平成26年度時点でも「M7を超えない」の反例である。\cite{jma_hyogo_history}
一つの活断層の平均活動間隔と全国の大地震の間隔を混ぜない。長期評価は過去の記録と推定幅を伴い、数千年の平均でも近い将来の発生を排除しない。原典の「最も活発でも約1000年」は母集団と期間の定義の不足を残す。\cite{herp_longterm_intervals}

\subsection{試験の読み方と次の単元への接続}
この章では生成場所、層の分類軸、数値の桁、断層の相対運動が判別点となる。頻度の最終評価は全年度の分類・本文確認後に行う。地震波の速度・震源過程は地震学、熱伝導と対流の原理は熱力学、岩石の実組成は岩石学へ接続する。各問題に残す解釈疑義は、この一般理論で公式以外の全肢を無条件に正しいとした記録ではない。
